% -----------------------------------------------------------------------------
% Results fragment, theme "misc": the small method figures. Four blocks, each tagged with the
% section it belongs to. Numbers: paper/results/scurve.json (scripts/paper/fig_scurve.R),
% eo_transforms.json (fig_eo_transforms.R), hs6_crosswalk.json (fig_hs6_crosswalk.R, which also
% writes tables/misc_hs6_counts.tex, Table tab:hs6_counts) and si_fixed_inputs.json
% (fig_si_fixed_inputs.R). Inputs: the final rerun of 12 October 2026 (data/processed), the HS6
% master data/reference/energy_hs6_master.csv after the October 2026 basket review and the
% Upstream baskets of 11 October, and the all-reporter Comtrade dyads of 2024.
% -----------------------------------------------------------------------------

% === section: normalization (Section 4, the S-curve transform) ===============
\paragraph{The S-curve in numbers.} With $\gamma=0.5$ the transform of
equation~\eqref{eq:scurve} maps the tenth and ninetieth percent ranks to 0.25 and 0.75 and the
quartiles to 0.366 and 0.634, so the middle half of a distribution occupies 0.27 of the unit
interval, against 0.5 under the plain percent rank. Its slope at the median is $\gamma=0.5$, half
that of the percent rank, and the slope rises without bound towards either end. Over a large
uniform grid of ranks the standard deviation of the transformed values tends to 0.19, against
0.29 for the percent rank and 0.38 for $\gamma=2$. In panel (b) of Figure~\ref{fig:scurve}, a
lognormal draw of 40 observations, min-max rescaling puts 29 of the 40 below 0.1 and spreads the
top four over 0.79 of the scale; the percent rank gives every group of four the same spread of
0.077; the S-curve gives the bottom four and the top four 0.22 each and the middle four 0.04.
The transform therefore keeps the order of the percent rank and its insensitivity to outliers,
and it reallocates resolution from the middle of the distribution to its ends. The small-sample
edge cases follow from the percent rank: a single observation scores 0, two score 0 and 1, a
fully tied set scores 0.5, and tied zeros below a positive value share 0.41.

\begin{figure}[tbp]
  \centering
  \includegraphics[width=\textwidth]{fig_scurve}
  \caption{The median-centred S-curve. Panel (a): $\Scurve(r)=r^{\gamma}/[r^{\gamma}+(1-r)^{\gamma}]$
  against the percent rank $r$ for $\gamma\in\{0.25,0.5,1,2\}$. The highlighted curve is
  $\gamma=0.5$, the value used for every component of the indices, and the marked points are
  $\Scurve(0.1)=0.25$, $\Scurve(0.5)=0.5$ and $\Scurve(0.9)=0.75$; $\gamma=1$ is the percent rank
  itself. Panel (b): 40 draws of a right-skewed variable (lognormal, log-mean 0, log-standard
  deviation 1, fixed seed) under three normalisations, with each observation joined across the
  rows: a min-max rescaling of the raw value, the percent rank $r$ of
  equation~\eqref{eq:pct_rank}, and the S-curved index of equation~\eqref{eq:scurve} computed by
  the pipeline's own function. The four lowest and four highest observations are highlighted, and
  the brackets give the span the S-curve assigns to each group. No pipeline data are used.
  Source: \texttt{scripts/paper/fig\_scurve.R}.}
  \label{fig:scurve}
\end{figure}

% === section: normalization (cardinal EO transforms; the float sits in Section 6) =========
\paragraph{The cardinal transformations of the opportunity pillar.} Three $\EO$ inputs are
cardinal scores rather than S-curved ranks (Figure~\ref{fig:eo_transforms}). The technology
readiness score rests on the 2026 IEA Clean Energy Technology Guide \citep{iea2026cleantechguide}:
of its 634 items, 253 map to a headline technology, and 9 of these to more than one
technology--stage cell. The Guide stops at TRL 9 and grades the items at that level by an
adoption stage, which the theme reads back onto an eleven-step scale; 98 items carry such a
stage, and 43 percent of the mapped items sit at TRL 9 or above. Because the level score is a bell
that peaks at TRL 6 and falls to zero at TRL 11, the 26 items that are commercial in many markets
score zero on level: the score rewards technologies that are proven but not yet mature, the
window in which entry is still open. The momentum score is concave in the TRL change over
2020--2025 and is capped at three steps; 54 percent of the items with a start-year reading did
not move. Where no item maps to a stage, the cell takes its technology's value: 10 of the 33
headline cells are scored this way, and the resulting index ranges from 0.31 (Geothermal
downstream) to 0.72 (Gas downstream), median 0.51. The overcapacity index $2^{-\rho}$ halves
with each unit of the capacity-to-demand ratio $\rho$. Of the four midstream technologies the
BloombergNEF supply-chain data cover, solar ($\rho=2.3$, index 0.20), batteries and wind ($\rho=3.0$,
index 0.13 each) sit on the steep part of the curve, and green hydrogen ($\rho=18.2$) is
effectively at zero. These scores take one value per technology--stage cell, so they move
pillar levels across cells but do not reorder economies within a cell.

\begin{figure}[tbp]
  \centering
  \includegraphics[width=\textwidth]{fig_eo_transforms}
  \caption{The three cardinal transformations of the Economic Opportunity pillar. Panel (a): the
  level score $b(x)$ of equation~\eqref{eq:eo_trl_bell} against the end-year technology readiness
  level $x$ (2025); circles sit on the curve at each level, with area proportional to the number
  of IEA Clean Energy Technology Guide items mapped to a headline technology at that level (253
  items). TRL 9 to 11 are the Guide's post-TRL-9 adoption stages: awaiting adoption (AA, 9),
  building momentum (BM, 10) and commercial in many markets (CM, 11). Panel (b): the momentum
  score $\psi(\Delta)$ against the change in TRL between 2020 and 2025, capped at three steps,
  with circles for the items that have a 2020 reading. Panel (c): the overcapacity index
  $2^{-\rho}$ of equation~\eqref{eq:eo_overcap} against the ratio $\rho$ of midstream
  manufacturing capacity to forecast demand, with the four technologies of the BloombergNEF
  supply-chain data at their values in the final build (green hydrogen, at $\rho=18.2$, lies off
  the scale). Parameters are the theme builders' own defaults. Source: \texttt{data/raw} (the
  Guide), \texttt{data/processed/technological\_readiness\_tbl.rds} and
  \texttt{overcapacity\_premium\_tbl.rds}; \texttt{scripts/paper/fig\_eo\_transforms.R}.}
  \label{fig:eo_transforms}
\end{figure}

% === section: data (Section 3, the HS6 crosswalk; replaces the pre-review table there) =====
\paragraph{The basket review.} Table~\ref{tab:hs6_counts} and Figure~\ref{fig:hs6_crosswalk}
describe the crosswalk after the October 2026 review. The review tested each code for dominant
use: a code stays in a technology's basket only if that technology is a main use of the product.
Ninety-two basket rows, touching 69 distinct codes, failed the test; 47 of these codes left every
basket and the others stay in another technology's basket. The rows are kept in the master with
the reason, so the review is reproducible and reversible. The review added 17 rows and moved six
from a midstream to an upstream basket, and the crosswalk went from 298 distinct codes in 490
basket rows to 259 codes in 415 rows: 242 of the current nomenclature and 17 legacy bridge codes,
49 upstream and 210 midstream. The mapping remains many-to-many: 105 codes (41 percent) belong to
more than one basket, 63 to two, 33 to three and 9, the transformer, switchgear and cable
headings, to four.

The exclusions matter most where general-purpose goods dominated a basket's trade. Measured on
the reconciled 2024 Comtrade matrix, the excluded codes carried 92 percent of the pre-review
coal midstream basket, 89 percent of the batteries upstream basket (copper ores alone, which now
belong to the grid's material basket, were USD 106 billion), 69 percent of gas midstream, 60
percent of oil midstream and 48 percent of batteries midstream, against 5 and 9 percent for the
solar and wind midstream baskets. Two changes redefine a basket. Electric Grid upstream held
only electrical energy (HS 2716.00, USD 61.9 billion in 2024), which made net power exporters
the grid's leading suppliers; it now holds seven copper and aluminium inputs, from ores and
bauxite to cathodes and unwrought aluminium, worth USD 321 billion. Switching and routing
apparatus (HS 8517.62), USD 219 billion and 32 percent of the pre-review Electric Grid midstream
basket, is telecommunication equipment and left it. Five upstream baskets are new: Electric
Vehicles (eight codes, the battery minerals and the rare earths of permanent magnets), Wind
(three, rare earths), Nuclear (two, uranium ores and natural uranium compounds), Solar (one,
silicon metal) and Semiconductors (gallium, germanium and indium, carried for description only).
The four that are scored trade little, from USD 3.0 billion (wind) to USD 21.3 billion (electric
vehicles) in 2024, and in uranium 26 percent of the value rests on exporters' declarations because
the importer did not report. Every scored upstream and midstream pair now has a basket
except those of Geothermal and Green Hydrogen.

\begin{figure}[tbp]
  \centering
  \includegraphics[width=\textwidth]{fig_hs6_crosswalk}
  \caption{The HS6 baskets after the October 2026 review. Panel (a): distinct HS6 codes in each
  technology--stage basket, split into codes of that basket only, current codes shared with at
  least one other basket (the many-to-many part of the mapping), legacy HS bridge codes, and codes
  added by the review (codes moved from a technology's midstream basket count as added to its
  upstream basket); left of zero, the codes the review excluded from the basket and the codes it
  moved to the upstream stage. Panel (b): share of the pre-review basket's 2024 trade value carried
  by the excluded codes; ``new'' marks a basket the review created. Trade is the reconciled
  all-reporter Comtrade matrix (the importer's declaration, else the exporter's declaration
  $\times1.06$), summed over exporter--importer pairs; a code in several baskets counts in each.
  Green hydrogen, geothermal and the downstream stage have no basket; semiconductors are carried
  in the master but not scored. Source: \texttt{data/reference/energy\_hs6\_master.csv},
  \texttt{data/raw/comtrade\_energy\_dyads\_all.csv}; \texttt{scripts/paper/fig\_hs6\_crosswalk.R}.}
  \label{fig:hs6_crosswalk}
\end{figure}

% Generated by scripts/paper/fig_hs6_crosswalk.R; do not edit by hand.
\begin{table}[tbp]
\centering
\begin{threeparttable}
\caption{HS6 baskets by technology and stage after the October 2026 basket review}
\label{tab:hs6_counts}
\footnotesize
\setlength{\tabcolsep}{3.5pt}
\begin{tabular}{@{}llrrrrrrrrr@{}}
\toprule
 & & \multicolumn{4}{c}{Codes in the basket} & \multicolumn{3}{c}{Review} & \multicolumn{2}{c}{Basket trade, 2024} \\
\cmidrule(lr){3-6}\cmidrule(lr){7-9}\cmidrule(l){10-11}
Technology & Stage & All & Legacy & Shared & Sub-sectors & Added & Moved & Excluded & USD bn & \% removed\\
\midrule
Batteries & Upstream & 12 & 1 & 5 & 2 & 0 & 0 & 20 & 21.2 & 88.8\\
 & Midstream & 34 & 0 & 2 & 9 & 0 & 0 & 11 & 202.7 & 48.4\\
Coal & Upstream & 6 & 0 & 0 & 2 & 0 & 0 & 0 & 192.5 & 0.0\\
 & Midstream & 15 & 1 & 12 & 3 & 0 & 0 & 8 & 20.5 & 91.9\\
Electric Grid & Upstream & 7 & 0 & 0 & 3 & 7 & 0 & 1 & 321.0 & 100.0\\
 & Midstream & 47 & 6 & 36 & 12 & 0 & 0 & 4 & 449.6 & 34.3\\
Electric Vehicles & Upstream & 8 & 0 & 8 & 8 & 8 & 0 & 0 & 21.3 & new\\
 & Midstream & 41 & 5 & 17 & 7 & 0 & 2 & 0 & 1,003.2 & 0.0\\
Gas & Upstream & 6 & 0 & 0 & 4 & 0 & 0 & 1 & 445.3 & 0.4\\
 & Midstream & 51 & 5 & 46 & 2 & 0 & 0 & 14 & 151.6 & 69.1\\
Nuclear & Upstream & 2 & 0 & 0 & 2 & 2 & 0 & 0 & 11.4 & new\\
 & Midstream & 18 & 1 & 12 & 3 & 0 & 1 & 0 & 24.6 & 0.0\\
Oil & Upstream & 10 & 1 & 0 & 5 & 0 & 0 & 0 & 2,211.0 & 0.0\\
 & Midstream & 43 & 4 & 35 & 7 & 0 & 0 & 14 & 127.1 & 59.8\\
Solar & Upstream & 1 & 0 & 0 & 1 & 1 & 0 & 0 & 3.8 & new\\
 & Midstream & 46 & 4 & 38 & 9 & 0 & 0 & 9 & 711.5 & 5.5\\
Wind & Upstream & 3 & 0 & 3 & 3 & 3 & 0 & 0 & 3.0 & new\\
 & Midstream & 44 & 3 & 39 & 8 & 0 & 2 & 10 & 409.5 & 9.4\\
Semiconductors\tnote{a} & Upstream & 2 & 1 & 0 & 1 & 2 & 0 & 0 & 1.1 & new\\
 & Midstream & 19 & 0 & 8 & 8 & 0 & 1 & 0 & 1,348.1 & 0.0\\
\midrule
Distinct codes & & 259 & 17 & 105 & 92 & 20 & 4 & 69 & & \\
\bottomrule
\end{tabular}
\begin{tablenotes}[flushleft]\footnotesize
\item \textit{Notes:} Distinct HS6 codes in each technology--stage basket of \texttt{data/reference/energy\_hs6\_master.csv} after the review. ``Legacy'' codes are bridge codes of older HS revisions whose successors are all in the basket; ``Shared'' counts the codes that also belong to at least one other basket. ``Excluded'' counts the codes the review removed from the basket because the technology is not a dominant use of the product, and ``Added'' the codes it added (the copper and aluminium inputs that replace electrical energy, HS 2716.00, in electric-grid Upstream, and the Upstream baskets of electric vehicles, solar, wind, nuclear and semiconductors, which had none); ``Moved'' counts the codes it moved from a technology's Midstream to its Upstream basket (rare-earth compounds and metals, natural uranium compounds, unwrought gallium, germanium and indium), which count as added there, since a code sits in one stage of a technology; the rows stay in the master with the reason. Basket trade is the 2024 value in the reconciled all-reporter Comtrade matrix (the importer's declaration, else the exporter's mirror declaration $\times$ 1.06), summed over exporter--importer pairs; ``\% removed'' is the share of the pre-review basket's value carried by the excluded codes (moved codes are not removals; ``new'' marks a basket the review created). A code in several baskets counts in each, so neither codes nor trade values add across rows. Green hydrogen, geothermal and the Downstream stage have no basket.
\item[a] Carried in the master; not a scored technology.
\end{tablenotes}
\end{threeparttable}
\end{table}


% === section: strategic (Section 8, the country-invariant inputs of the SI) ================
\paragraph{The fixed inputs of the Strategic Index.} Two of the five terms of the Strategic Index
take the same value for every country, and they set a sector prior
(Figure~\ref{fig:si_fixed_inputs}, panel b). The stage term gives midstream sectors 0.05 more
than downstream ones; the technology term, $0.1\,g_t$ with $g_t$ the clean-technology weight
built from IPCC life-cycle emissions, separates oil ($g_t=0$) from wind ($g_t=0.999$) by 0.10.
The prior thus runs from 0.025 (Oil downstream) to 0.175 (Wind midstream), out of a possible range
of the index of 0.025 to 0.975. The prior is a stated preference for clean, midstream
activity, not a measurement, and its weight is visible when the country-varying terms are flat.
For a country whose only rankable sectors are the coal, gas, nuclear and oil midstream pairs, the
prior alone puts nuclear 0.058 ahead of gas, 0.097 ahead of coal and 0.100 ahead of oil; with
midstream nuclear security compressed around its median, this lead, rather than a floor in
security, is what places nuclear midstream at the top of such countries' Strategic Index
(Section~\ref{sec:strategic}). The Climate Action Tracker component of the
Policy index (panel a) uses the same weight $g_t$: a country's rating $\kappa_i$ sets the slope
$2\kappa_i-1$, so a rating of ``insufficient'' or better favours clean technologies and a rating of
``critically insufficient'' favours fossil ones. In the final build CAT rates 65 countries, 26 of
them EU members through the EU's rating, and 46 of the 57 coalition members; 31 members are
rated ``insufficient'', 7 ``critically insufficient'', 4 ``highly insufficient'' and 4 ``almost
sufficient''.

\begin{figure}[tbp]
  \centering
  \includegraphics[width=\textwidth]{fig_si_fixed_inputs}
  \caption{The country-invariant inputs of the Strategic Index. Panel (a): the Climate Action
  Tracker component of the Policy index, $x^{\mathrm{CAT}}_{its}=\kappa_ig_t+(1-\kappa_i)(1-g_t)$,
  against the technology weight $g_t$ (IPCC life-cycle emissions, 1 = cleanest) for the four CAT
  ratings, with the eleven technologies marked on each line; the slope is $2\kappa_i-1$. The
  legend counts the countries CAT rates in each category in the final build (65 in all, 26 of
  them EU members through the EU's rating) and the coalition members among them (46 of 57).
  Panel (b): the sector prior $0.1\lambda_s+0.1g_t$ that the stage weights $\lambda_s$ (upstream
  0.5, midstream 0.75, downstream 0.25) and the technology weights add to the Strategic Index of
  every country; Geothermal midstream is not a reported sector. Source: the pipeline's own
  functions and \texttt{data/raw/ipcc\_ghg\_intensity.csv}, with the CAT ratings of
  \texttt{policy\_component\_tbl} in \texttt{data/processed/outputs/index\_outputs.rds};
  \texttt{scripts/paper/fig\_si\_fixed\_inputs.R}.}
  \label{fig:si_fixed_inputs}
\end{figure}
